\section{Binary Expand--Accumulate}
\label{sec:binary-ea}

We first keep the binary EA construction of~\cite{EA} unchanged.  The only
change is the analysis.  For each message weight, the Bernoulli expander gives
independent intermediate bits.  The accumulator then has an exact two-state
description.

\subsection{Exact first moment}

Fix a message $\x\in\F_2^k$ of weight $r$ and sample
$\B\gets\mathcal B^{\mathrm{Ber}}_{k,n,\rho}$.  Each coordinate of $\x\B$ is
the parity of $r$ independent Bernoulli variables.  Its probability of being
one is
\begin{equation}
 a_r:=\frac{1-(1-2\rho)^r}{2}.
 \label{eq:ea-activation}
\end{equation}
Different coordinates use independent columns of $\B$.  Thus the entries of
$\x\B$ are independent $\Ber(a_r)$ variables.

Let $Y$ mark a nonzero accumulator output.  Rows and columns of
\begin{equation}
 \mathbf T_a(Y):=
 \begin{pmatrix}
  1-a & aY\\
  a   & (1-a)Y
 \end{pmatrix}
 \label{eq:ea-transfer}
\end{equation}
index the current and next accumulator states in the order $0,1$.  The factor
$Y$ marks the next state because that state is also the current output symbol.
Let $\e_0:=(1,0)^{\mathsf T}$.

\begin{theorem}[Exact Bernoulli EA first moment]
\label{thm:ea-exact-first-moment}
For $\G=\B\A$ and every $0\leq L\leq n$, define
\begin{equation}
 \mu_{\EA}(L):=
 \sum_{r=1}^{k}\binom{k}{r}
 \sum_{h=0}^{L}[Y^h]\,
 \e_0^{\mathsf T}\mathbf T_{a_r}(Y)^n\one.
 \label{eq:ea-first-moment}
\end{equation}
Then
\begin{equation}
 \Pr_{\B}[\Bad_L(\B\A)]\leq\mu_{\EA}(L).
 \label{eq:ea-first-moment-failure}
\end{equation}
\end{theorem}

\begin{proof}
For a fixed weight-$r$ message, the matrix product in
\eqref{eq:ea-first-moment} is the probability generating function of its
output weight.  Coefficient extraction sums its probability of output weight
at most $L$.  There are $\binom{k}{r}$ binary messages of weight $r$.
Equation~\eqref{eq:first-moment} gives the probability bound.
\end{proof}

The theorem retains the exact accumulator process.  The next bound replaces
coefficient extraction by a one-dimensional optimization without changing
the transfer matrix.

For $a\in[0,1/2]$ and $z\in(0,1]$, define
\begin{align}
 \lambda(a,z)
 &:=\frac{(1-a)(1+z)+
 \sqrt{(1-a)^2(1-z)^2+4a^2z}}{2},
 \label{eq:ea-spectral-radius}\\
 J_\delta(a)&:=
 \sup_{0<z\leq1}\{\delta\ln z-\ln\lambda(a,z)\}.
 \label{eq:ea-rate-function}
\end{align}
The quantity $\lambda(a,z)$ is the spectral radius of
$\mathbf T_a(z)$, meaning the largest absolute value of its eigenvalues.

\begin{lemma}[Accumulator lower tail]
\label{lem:ea-lower-tail}
Let $H$ be the output weight of the binary accumulator on $n$ independent
$\Ber(a)$ inputs.  For $a\in(0,1/2]$ and $\delta\in(0,1/2)$,
\begin{equation}
 \Pr[H\leq\lfloor\delta n\rfloor]
 \leq\sqrt2\exp(-nJ_\delta(a)).
 \label{eq:ea-lower-tail}
\end{equation}
Moreover,
\begin{align}
 \liminf_{a\downarrow0}\frac{J_\delta(a)}a
 &\geq I_\delta:=1-2\sqrt{\delta(1-\delta)},
 \label{eq:ea-sparse-endpoint}\\
 J_\delta(1/2)&=\ln2-h(\delta).
 \label{eq:ea-dense-endpoint}
\end{align}
The function $J_\delta$ is continuous and positive on $(0,1/2]$.
\end{lemma}

\begin{proof}
Set $D_z:=\diag(1,\sqrt z)$.  The similar matrix
\[
 D_z\mathbf T_a(z)D_z^{-1}
 =\begin{pmatrix}1-a&a\sqrt z\\a\sqrt z&(1-a)z\end{pmatrix}
\]
is symmetric and positive semidefinite.  Its largest eigenvalue is
$\lambda(a,z)$.  Hence
\[
 \E[z^H]
 =\e_0^{\mathsf T}\mathbf T_a(z)^n\one
 \leq\sqrt{1+z}\,\lambda(a,z)^n.
\]
Apply \eqref{eq:positive-marker-tail} and optimize over $z$.

For the sparse endpoint, set $z=e^{-\theta a}$ and expand at $a=0$.
The coefficient of $a$ is optimized at
$\theta=(1-2\delta)/\sqrt{\delta(1-\delta)}$, which gives
\eqref{eq:ea-sparse-endpoint}.  At $a=1/2$, one has
$\lambda(1/2,z)=(1+z)/2$.  The choice $z=\delta/(1-\delta)$ gives
\eqref{eq:ea-dense-endpoint}.  For $a>0$, the derivative of
$\ln\lambda(a,z)$ with respect to $\ln z$ at $z=1$ is the stationary
nonzero-state probability $1/2>\delta$.  A marker slightly below one
therefore gives a positive exponent.  Also, $\lambda(a,z)\geq1-a\geq1/2$,
so the defining objective is at most $\delta\ln z+\ln2$.
This tends uniformly to $-\infty$ as $z\downarrow0$.  Restricting the
optimization to a compact interval of positive markers proves continuity.
\end{proof}

\subsection{Asymptotic distance}

The endpoint formulas separate the sparse and dense message regimes.  Sparse
messages determine the required expander density.  Dense messages give the GV
condition.

\begin{theorem}[Binary EA distance]
\label{thm:ea-distance}
Let $k=k(n)$ satisfy $k/n\to R\in(0,1)$.  Sample
$\B\gets\mathcal B^{\mathrm{Ber}}_{k,n,\rho_n}$ with
\begin{equation}
 \rho_n=C\frac{\ln n}{n}
 \label{eq:ea-density}
\end{equation}
for a constant $C>0$.  Fix $\delta\in(0,1/2)$.  If
\begin{align}
 C I_\delta&>1,
 \label{eq:ea-sparse-condition}\\
 h(\delta)&<(1-R)\ln2,
 \label{eq:ea-gv-condition}
\end{align}
then
\begin{equation}
 \Pr_{\B}\!\left[d_{\min}(\B\A)\leq\lfloor\delta n\rfloor\right]
 =o(1).
 \label{eq:ea-distance-conclusion}
\end{equation}
\end{theorem}

\begin{proof}
We apply Lemma~\ref{lem:ea-lower-tail} to
Theorem~\ref{thm:ea-exact-first-moment} and split the message weights into
three ranges.

Choose constants $\varepsilon,\tau>0$ and $\gamma>1$ such that
$C(1-\tau)(I_\delta-\varepsilon)>\gamma$.
Decrease $\tau$ until $J_\delta(a)\geq a(I_\delta-\varepsilon)$ for
$0<a\leq\tau$, as permitted by \eqref{eq:ea-sparse-endpoint}.
For all sufficiently large $n$ and every $r\rho_n\leq\tau$,
\[
 r\rho_n(1-r\rho_n)
 \leq\frac{1-e^{-2r\rho_n}}2\leq a_r\leq r\rho_n.
\]
Lemma~\ref{lem:ea-lower-tail} therefore bounds the weight-$r$ tail by
$\sqrt2\,n^{-\gamma r}$.  Since $\binom{k}{r}\leq k^r$ and $k=O(n)$,
summing this geometric bound over the sparse range gives $o(1)$.

For $r\rho_n>\tau$, one has $a_r\geq(1-e^{-2\tau})/2$.  Positivity
and continuity of $J_\delta$ give a uniform exponential tail.  The number of
messages with $r\leq\eta n$ has exponent
$R h(\eta/R)+o(1)$.  Choose $\eta>0$ small enough that this counting exponent
is below the tail exponent.  The intermediate contribution is then
exponentially small.

For $r>\eta n$, equation~\eqref{eq:ea-activation} gives
$a_r=1/2-o(1)$ uniformly.  There are at most $2^k$ messages.  The dense
endpoint bounds their total contribution by
\[
 \exp\!\left(n\bigl(R\ln2+h(\delta)-\ln2+o(1)\bigr)\right),
\]
which vanishes by \eqref{eq:ea-gv-condition}.  The three ranges cover every
nonzero message weight.
\end{proof}

Condition~\eqref{eq:ea-gv-condition} is exactly the binary GV condition.
The accumulator therefore loses no asymptotic exponent for dense messages.
The sparse condition is an additional cost of the sparse expander.

At $\delta=0.05$, the sparse condition requires $C>1.773$.  The analysis in
\cite[Theorem~3.10]{EA} requires $C>4.938$.  At rate $R=1/5$, its rate
condition permits $\delta<0.1127$, while
\eqref{eq:ea-gv-condition} permits $\delta<0.2430$.  These are asymptotic
comparisons; they do not state a finite failure probability.

\subsection{Finite same-ensemble certificates}

The verifier described in Section~\ref{sec:certificates} evaluates
\eqref{eq:ea-first-moment} over every message weight.
Its block bound uses the following monotonicity argument.
Write $S_a(z):=D_z\mathbf T_a(z)D_z^{-1}$ for the symmetric matrix in
Lemma~\ref{lem:ea-lower-tail}.  Its derivative is
\[
 \frac{\partial S_a(z)}{\partial a}
 =-\begin{pmatrix}1\\-\sqrt z\end{pmatrix}
   \begin{pmatrix}1&-\sqrt z\end{pmatrix}.
\]
Every quadratic form of this derivative is nonpositive.  Thus the largest
eigenvalue $\lambda(a,z)$ is nonincreasing in $a$.
For $0<\rho<1/2$, the activation probability $a_r$ is nondecreasing in $r$.
Consequently, for any integer interval $[u,v]\subseteq[1,k/2]$ and
$0<z\leq1$,
\begin{equation}
 \sum_{r=u}^{v}\binom{k}{r}\Pr[H_r\leq L]
 \leq(v-u+1)e^{k h(v/k)}\sqrt{1+z}\,z^{-L}\lambda(a_u,z)^n,
 \label{eq:ea-finite-block}
\end{equation}
where $H_r$ is the accumulator-output weight for a fixed weight-$r$ message.
Here $\binom{k}{r}\leq e^{k h(r/k)}$, and $h$ is nondecreasing on $[0,1/2]$.
For a block above $k/2$, replace the count factor
$(v-u+1)e^{k h(v/k)}$ by $2^k$.

\begin{theorem}[Finite Bernoulli EA certificates]
\label{thm:ea-finite}
For each row of Table~\ref{tab:ea-finite}, set $k:=2^s$, $n:=5k$,
$\rho:=w/n$, and $L:=n/20$.  If
$\B\gets\mathcal B^{\mathrm{Ber}}_{k,n,\rho}$, then
\begin{equation}
 \Pr_{\B}[\Bad_L(\B\A)]\leq U<2^{-20}.
\end{equation}
\end{theorem}

\begin{table}[t]
\centering
\caption{Certified binary Bernoulli EA bounds at rate $1/5$ and relative
distance $0.05$.  The parameter $w$ is the expected expander row weight.}
\label{tab:ea-finite}
\begin{tabular}{@{}crrr@{}}
\toprule
$s$ & $w$ & $U$ & $-\log_2U$ \\
\midrule
$20$ & $50$ & $6.891818499\cdot10^{-7}$ & $20.468611$ \\
$25$ & $56$ & $7.475946766\cdot10^{-7}$ & $20.351240$ \\
$30$ & $62$ & $8.107513660\cdot10^{-7}$ & $20.234237$ \\
\bottomrule
\end{tabular}
\end{table}

\begin{proof}
For each row, the verifier evaluates weights $1$ through $128$ with the exact
transfer matrix.  It partitions the remaining weights into contiguous
intervals and applies \eqref{eq:ea-finite-block} or its $2^k$ count variant.
Their union is $[129,k]$.  Outward-rounded interval arithmetic
evaluates every stored marker and the final sum.  The verified upper endpoint
is the displayed value $U$.  Theorem~\ref{thm:ea-exact-first-moment} converts
that sum into the claimed probability bound.
\end{proof}

At $s=20$, the original $C=3$ profile has expected row weight $46.42$ but an
extrapolated failure bound $0.0157$ at distance $0.05$~\cite{EA}.  The first
row of Table~\ref{tab:ea-finite} is slightly denser and proves a complete
$2^{-20}$ statement.  At $s=25$ and $s=30$, the certified profiles are both
sparser and have much smaller failure bounds than the corresponding published
profiles.
